\section{Proofs for the Parameter-free Case (Section \ref{sec:parameter-free})}
\label{app:parameter-free}

% Previously, we always bound $\|w_t-x^*\|\leq D$, but the following lemma gives a finer bound on the inner product, which preserves the quantity $\|w_t\|$ and $\|x^*\|$. This allows us to come up with an adaptive bound by regularizing $\ell_t(w)$ by $\|w\|^2$.
% In this subsection, we only consider the specific case in 2-norm, and we make the following definition:
\begin{definition}
A random vector $X\in\RR^d$ is said to be $\sigma$-norm-sub-Gaussian, denoted by $\nSG(\sigma)$ if
\begin{align*}
    \P\{\|X-\E[X]\|_2 \geq \epsilon\} \leq 2\exp\left(-\frac{\epsilon^2}{2\sigma^2}\right), \ \forall \epsilon.
\end{align*}
\end{definition}

We will rely on the following concentration bound on norm-sub-Gaussian random vectors.

\begin{lemma}[Lemma 1, \cite{jin2019short}]
\label{lem:normSG}
There exists a universal $C$ such that (i) if $\|X\|\leq \sigma$, then $X$ is $\nSG(C\sigma)$ nad (ii) if $X$ is $\sigma$-sub-Gaussian, then $X$ is $\nSG(C\sqrt{d}\sigma)$. 
\end{lemma}

\begin{lemma}[Corollary 8, \cite{jin2019short}]
\label{lem:nSG-concentration}
There exists a universal constant $C$ such that if $X_i|X_{1:i-1}$ is mean-zero $\nSG(\sigma_i)$ for all $X_1,\ldots,X_t$, then for any fixed $\delta>0$ and $B>b>0$ such that $b<\sum_{i=1}^t \sigma_i^2 \leq B$ almost surely, with probability at least $1-\delta$,
\begin{align*}
    \left\|\sum_{i=1}^t X_i\right\|_2 \leq C\sqrt{\sum_{i=1}^t \sigma_i^2\left(\log\frac{2d}{\delta}+\log\log\frac{B}{b} \right)}.
\end{align*}
\end{lemma}

Recall that $\weight_t\nabla f(x_t)-g_t=\sum_{i=1}^t X_i$ (Lemma \ref{lem:variance-expectation-bound}), where
\[
X_i = [\weight_i\nabla f(x_i)-\weight_{i-1}\nabla f(x_{i-1})]-[\weight_i\nabla\ell(x_i,z_i)-\weight_{i-1}\nabla\ell(x_{i-1},z_i)];
\]
and $\gamma_t = \sum_{i\in I_t}R_i$, where $R_i=\sigma_i\tilde R_i$ and $R_i\sim \calD$ i.i.d. We have the following lemma:

\begin{lemma}
\label{lem:X-R-nSG}
Suppose Assumption \ref{as:bounded-domain} - \ref{as:smooth} hold w.r.t. the 2-norm, and suppose $\calD$ is a $(V,\alpha)$-RDP distribution and is $\sigma_\calD$-sub-Gaussian, i.e.,
\begin{align*}
    \P\{\sup_{\|a\|=1} \langle X, a\rangle \geq \epsilon \} \leq \exp\left(- \frac{\epsilon^2}{2\sigma_\calD^2}\right).
\end{align*}
Also set $\weight_t=t^k$. Then there exists a universal constant $C$ such that $X_i|X_{1:i-1}$ are mean-zero $\nSG(\sigma_{X_i})$ and $R_i|R_{1:i-1}$ are mean-zero $\nSG(\sigma_{R_i})$ for all $i$, where
\begin{align*}
    % & \sigma_{X_i} = 2C|\weight_i-\weight_{i-1}|G+2C(\weight_i^2/\weight_{1:i})H\|w_i-x_{i-1}\|_2, \\
    & \sigma_{X_i} = 2C(k+1)(G+H\|w_i-x_{i-1}\|_2)i^{k-1}, \\
    & \sigma_{R_i} = C\sqrt{d}\sigma_\calD \sigma_i.
\end{align*}
\end{lemma}

\begin{proof}
Since we assume $\E[\nabla\ell(x,z)]=\nabla f(x)$ for all $x$, $\E[X_i|X_{1:i-1}]=0$. Also, since $\calD$ is a $(V,\alpha)$-RDP distribution (Definition \ref{def:RDP-dist}) and $R_i=\sigma_i\tilde R_i$'s are independent, $\E[R_i|R_{1:i-1}]=\E[\tilde R_i]=0$. 

For the second part, in Lemma \ref{lem:variance-expectation-bound} we proved that
\begin{align*}
    \|\delta_i\|_2 = \|\weight_i\nabla\ell(x_i,z_i)-\weight_{i-1}\nabla\ell(x_{i-1},z_i)\|_2 \leq (k+1)(G+H\|w_i-x_{i-1}\|_2)i^{k-1}.
\end{align*}
The same bound holds for $\weight_i\nabla f(x_i)-\weight_{i-1}\nabla f(x_{i-1})$ following the same argument. Therefore,
\[
% \|X_i\|_2 \leq 2|\weight_i-\weight_{i-1}|G+2(\weight_i^2/\weight_{1:i})H\|w_i-x_{i-1}\|_2.
\|X_i\|_2 \leq 2(k+1)(G+H\|w_i-x_{i-1}\|_2)i^{k-1}.
\]
Moreover, since we assume $\tilde R_i\sim \calD$ is $\sigma_\calD$-sub-Gaussian, $R_i=\sigma_i\tilde R_i$ is $\sigma_i\sigma_\calD$-sub-Gaussian.
Hence, by Lemma \ref{lem:normSG}, $X_i|X_{1:i-1}$ and $R_i|R_{1:i-1}$ are norm-sub-Gaussian.
\end{proof}

% \highlight{To be more careful in the proof, we may need to say $\F_t=\{(w_i,x_i)\}_{i=1}^{t+1}\cup\{(g_i,\gamma_i)\}_1^t$ is a filtration of $X_t$, then $X_i|\F_{t-1}$ is conditionally $\nSG(...\|w_i-x_{i-1}\|)$.}

\begin{theorem}
\label{thm:DP-OTB:parameter-free-2-norm-step-1}
Suppose w.r.t. 2-norm, $W$ is bounded by $D$ and $\ell$ is $G$-Lipschitz and $H$-smooth. Suppose $\calD$ is $(V,\alpha)$-RDP distribution and is $\sigma_\calD$-sub-Gaussian. If we set $\weight_t=t^3$ (i.e. $k=3$) and set $\sigma_t^2$ as defined in \eqref{eq:DP-OTB:RDP-noise}, then with probability at least $1-\delta$,
\begin{align*}
    f(x_T)-f(x^*) \leq \frac{4}{T^4}\left( \regret_T(x^*) + \sum_{t=1}^T \xi_t(\|w_t\|_2+\|x^*\|_2) + \nu_t(\|w_t\|_2^2+\|x^*\|_2^2) \right).
\end{align*}
where $C$ is a universal constant, $A=8\sqrt{2}C^2, A'=8\sqrt{d}\sigma_\calD C^2, \kappa=1+DH/G$, 
% $\Phi=\sqrt{\log(16dT\log(2\kappa T)/\delta)}$, 
and
\begin{align*}
    &\xi_t = AG\Phi t^{5/2} + A'(G+DH)\frac{\Phi\log_2(2T)t^2}{\rho} , \
    \nu_t = 28AH\Phi t^{5/2}, \ 
    \Phi = \sqrt{\log\frac{20dT\log(2\kappa T)}{\delta}}.
\end{align*}
\end{theorem}

\begin{proof}
We start with Eq. \eqref{eq:basic-decomposition}:
\begin{align}
    \weight_{1:T}(f(x_T)-f(x^*))
    &\leq R_T(x^*) + \sum_{t=1}^T\langle \weight_t\nabla f(x_t)-g_t-\gamma_t,w_t-x^*\rangle \notag\\
    % \intertext{By Fenchel-Young's inequality and triangle inequality,}
    &\leq R_T(x^*) + \sum_{t=1}^T \left(\|\weight_t\nabla f(x_t)-g_t\|_2+\|\gamma_t\|_2\right) (\|w_t-x^*\|_2). \label{eq:parameter-free-base}
\end{align}

\textbf{Step 1.}
By Lemma \ref{lem:nSG-concentration} and \ref{lem:X-R-nSG}, for each $t$, with probability $1-\delta/2T$,
\begin{align*}
    \|\weight_t\nabla f(x_t)-g_t\|_2 
    \leq C\sqrt{\sum_{i=1}^t \sigma_{X_i}^2 \left( \log\frac{4dT}{\delta} + \log\log\frac{B}{b}\right)}.
\end{align*}
Since we choose $\weight_t=t^3$ (i.e., $k=3$),
\begin{align*}
    \sigma_{X_i} = 8Ci^2(G+H\|w_i-x_{i-1}\|_2).
\end{align*}
% Note that for all $t$, $\sum_{i=1}^t \sigma_{X_i}^2 \leq B:=64C^2(G+DH)^2T^5$, and $\sum_{i=t}^t\sigma_{X_i}^2\geq b:=\sigma_{X_1}^2=4C^2(G+H\|w_1\|_2)^2$. 
Next, we can bound $\sum_{i=1}^t\sigma_{X_i}^2$ as follows: for all $t$,
\begin{align*}
    \textstyle \sum_{i=1}^t \sigma_{X_i}^2 &\leq B := 64C^2(G+DH)^2T^5, \\
    \textstyle \sum_{i=t}^t \sigma_{X_i}^2 &\geq b := \sigma_{X_1}^2 = 64C^2(G+H\|w_1\|_2)^2.
\end{align*}
Recall that $\kappa = 1+DH/G$, so
\begin{align*}
    \frac{B}{b} = \frac{(G+DH)^2T^5}{(G+H\|w_1\|_2)^2} \leq (\kappa T)^5.
\end{align*}
Also recall that $\Phi=\sqrt{\log(20dT\log(2\kappa T)/\delta)}$, so
\begin{align*}
    \sqrt{\log\frac{4dT}{\delta} + \log\log \frac{B}{b}} \leq \sqrt{\log \frac{20dT\log(\kappa T)}{\delta}} \leq \Phi.
\end{align*}
Therefore, with probability at least $1-\delta/2T$,
\begin{align}
    \|\weight_t \nabla f(x_t) - g_t\|_2 
    &\leq C\Phi \sqrt{\sum_{i=1}^t [8Ci^2(G+H\|w_i-x_{i-1}\|_2)]^2} \notag\\
    &\leq 8C^2\Phi \sqrt{\sum_{i=1}^t i^4(G+H\|w_i-x_{i-1}\|_2)^2}.
    \label{eq:variance-concentration}
\end{align}
By union bound, with probability at least $1-\delta/2$,
\begin{align}
    &\sum_{t=1}^T \|\weight_t\nabla f(x_t)-g_t\|_2\|w_t-x^*\|_2 \notag\\
    \leq & 8C^2\Phi \sum_{t=1}^T \sqrt{\sum_{i=1}^t i^4(G+H\|w_i-x_{i-1}\|_2)^2} \|w_t-x^*\|_2 \notag
    \intertext{We use the identity $(a+b)^2\leq 2a^2+2b^2$ and $\sqrt{a+b}\leq \sqrt{a}+\sqrt{b}$:}
    \leq & 8C^2\Phi \sum_{t=1}^T \left(\sqrt{\sum_{i=1}^t 2G^2i^4} + \sqrt{\sum_{i=1}^t 2H^2\|w_i-x_{i-1}\|_2^2i^4}\right) \|w_t-x^*\|_2 \label{eq:step-1}
\end{align}
\textbf{1.1.} We bound these two sums separately. For the first sum,
recall that $A=8\sqrt{2}C^2$, so
\begin{align*}
    8C^2\sum_{t=1}^T \sqrt{\sum_{i=1}^t 2G^2i^4}\|w_t-x^*\|_2 
    \leq AG\sum_{t=1}^T t^{5/2}(\|w_t\|_2+\|x^*\|_2).
\end{align*}

\textbf{1.2.}
For the second sum, we apply Young's inequality ($ab\leq \frac{1}{2\lambda}a^2+\frac{\lambda}{2}b^2$) for each $t$:
\begin{align}
    &8C^2\sum_{t=1}^T\sqrt{\sum_{i=1}^t2H^2\|w_i-x_{i-1}\|_2^2i^4}\|w_t-x^*\|_2 \notag\\
    \leq & AH\sum_{t=1}^T\frac{1}{2\lambda_t}\sum_{i=1}^t (\|w_i-x_{i-1}\|_2^2i^4) + \frac{\lambda_t}{2}\|w_t-x^*\|_2^2 \notag
    % line 1
    \intertext{We first bound $\|w_i-x_{i-1}\|_2^2 \leq 2\|w_i\|_2^2+2\|x_{i-1}\|_2^2$. Recall that $x_0=0$ and for $i\geq 2$, $x_{i-1}=\sum_{j=1}^{i-1} \frac{\weight_j}{\weight_{1:{i-1}}}w_j$, so $\|x_{i-1}\|_2^2\leq \sum_{j=1}^{i-1}\frac{\weight_j}{\weight_{1:i-1}}\|w_j\|_2^2$ by convexity. Consequently,}
    \leq & AH\sum_{t=1}^T \left( \frac{1}{\lambda_t}\sum_{i=1}^t i^4\|w_i\|_2^2 + \frac{1}{\lambda_t}\sum_{i=2}^t i^4\sum_{j=1}^{i-1} \frac{\weight_j\|w_j\|_2^2}{\weight_{1:i-1}} + \lambda_t(\|w_t\|_2^2+\|x^*\|_2^2) \right). \label{eq:step-1.2}
\end{align}
We define $\lambda_t = ct^{5/2}$ for some constant $c$ to be determined later, and we apply change of summation on the first two sums:
\begin{lemma}
\label{lem:change-sum}
For any sequence $a_i,b_j,c_k$,
\begin{align*}
    \sum_{i=1}^N a_i \sum_{j=1}^i b_j = \sum_{i=1}^N b_i \sum_{j=i}^N a_j,
    \quad \textit{and} \quad 
    \sum_{i=1}^N a_i \sum_{j=1}^i b_j \sum_{k=1}^j c_k = \sum_{i=1}^N c_i \sum_{j=i}^N a_j \sum_{k=i}^j b_k.
\end{align*}
\end{lemma}

\textbf{1.2.1.}  For the first summation,
\begin{align*}
    \sum_{t=1}^T \frac{1}{\lambda_t} \sum_{i=1}^t i^4\|w_i\|_2^2
    &= \sum_{t=1}^T\sum_{i=t}^T \frac{1}{\lambda_i} t^4\|w_t\|_2^2
    \intertext{For decreasing function $f$, $\sum_{i=t+1}^T f(i)\leq \int_t^T f(x)\,dx$, then:}
    &\leq \sum_{t=1}^T \left( \frac{1}{ct^{5/2}} + \int_t^\infty \frac{1}{cx^{5/2}}\,dx \right)t^4\|w_t\|_2^2
    % \leq \sum_{t=1}^T \left(\frac{2}{3c}t^{5/2} + \frac{1}{c}t^{3/2}\right) \|w_t\|_2^2. 
    \leq \sum_{t=1}^T \frac{5}{3c}t^{5/2}\|w_t\|_2^2.
\end{align*}
% In general, we define $\lambda_t\propto t^{k-\frac{1}{2}}$. However, for simplicity, we consider the case when $k=3$, and define $\lambda_t = ct^{k-\frac{1}{2}}$. 
% We then use the identity that $\sum_{i=t}^T f(i) \leq f(t)+\int_t^T f(x) dx$ for decreasing $f$:
% \begin{align*}
%     \sum_{i=t}^T \frac{c}{\lambda_i}
%     = \sum_{i=t}^T \frac{1}{i^{5/2}} \leq \frac{1}{t^{5/2}} + \int_t^\infty \frac{1}{x^{5/2}}\, dx = \frac{1}{t^{5/2}} + \frac{2}{3t^{3/2}} \leq \frac{5}{3t^{3/2}}.
% \end{align*}
% \begin{align*}
%     \sum_{i=t}^T \frac{1}{\lambda_i}\sqrt{\log \frac{2d}{\delta_i}}
%     &= \sum_{i=t}^T \frac{\sqrt{\log(\pi^2di^2/3\delta)}}{i^{5/2}}
%     \intertext{The summand is decreasing for all $i\geq \sqrt{\frac{3\delta\exp(2/5)}{\pi^2d}}$, which is less than $1$ for all $\delta\leq 1, d\geq 1$. We then use the identity that $\sum_{i=t}^T f(i) \leq f(t)+\int_t^T f(x) dx$ for decreasing $f$:}
%     &\leq \frac{\sqrt{\log(\pi^2dt^2/3\delta)}}{t^{5/2}} + \int_t^\infty \frac{\sqrt{\log(\pi^2dx^2/3\delta)}}{x^{5/2}} \, dx
%     \intertext{This integral is hard to evaluate. Alternatively, observe that $\pi^2d/3\delta \geq e$, so $\log(\pi^2di^2/3\delta)\geq 1$ and $\sqrt{\log(\pi^2di^2/3\delta)}\leq \log(\pi^2di^2/3\delta)$ for all $i\geq 1$. Hence:}
%     &\leq \frac{\sqrt{\log(\pi^2dt^2/3\delta)}}{t^{5/2}} + \int_t^\infty \frac{\log(\pi^2dx^2/3\delta)}{x^{5/2}} \, dx \\
%     &= \frac{\sqrt{\log(\pi^2dt^2/3\delta)}}{t^{5/2}} + \frac{8}{9t^{3/2}} + \frac{2\log(\pi^2dt^2/3\delta)}{t^{3/2}} \\
%     &\leq \frac{4\log(\pi^2dt^2/3\delta)}{t^{3/2}}.
% \end{align*}
% \highlight{Another way to evaluate this integral is by bounding $\sqrt{\log(\pi^2di^2/3\delta)}\leq \sqrt{\log(\pi^2dT^2/3\delta}$. This would preserve the order $\sqrt{\log}$, but it's no longer any-time and requires the knowledge of $T$, so I personally prefer the current bound.}
% Therefore, the first term is bounded by:
% \begin{align*}
%     \sum_{t=1}^T\sum_{i=t}^T \frac{1}{\lambda_i} 2\|w_t\|_2^2t^{2(k-1)}
%     % \leq \sum_{t=1}^T 8\log\left(\frac{\pi^2dt^2}{3\delta}\right)t^{5/2}\|w_t\|_2^2.
%     \leq \sum_{t=1}^T \frac{10}{3c}t^{5/2}\|w_t\|_2^2.
% \end{align*}
\textbf{1.2.2.} For the second term, by Proposition \ref{prop:alphas}, $\weight_{1:i-1}\geq (i-1)^4/4$, so
\begin{align*}
    \sum_{t=1}^T \frac{1}{\lambda_t} \sum_{i=2}^t i^4\sum_{j=1}^{i-1} \frac{\weight_j\|w_j\|_2^2}{\weight_{1:i-1}} 
    &\leq \sum_{t=1}^T \frac{1}{ct^{5/2}}\sum_{i=2}^{t} \frac{4i^4}{(i-1)^4} \sum_{j=1}^{i} j^3\|w_j\|_2^2 
    \intertext{For all $i\geq 2$, we can bound $i/(i-1)\leq 2$. We then apply change of summation, which gives:}
    &\leq \sum_{t=1}^T\sum_{i=t}^T\frac{1}{ci^{5/2}}\sum_{j=t}^i 64t^3\|w_t\|_2^2
    \leq \sum_{t=1}^T \frac{192}{c} t^{5/2}\|w_t\|_2^2.
\end{align*}
The last inequality is again derived from the integral bound:
\begin{align*}
    \sum_{i=t}^T \frac{1}{i^{5/2}}\sum_{j=t}^i 1 \leq \sum_{i=t}^T \frac{1}{i^{3/2}}
    \leq \frac{1}{t^{3/2}} + \int_t^\infty \frac{1}{x^{3/2}}\,dx
    \leq \frac{3}{t^{1/2}}.
\end{align*}
% We then bound the inner summation with the same argument using the integral inequality. By Proposition \ref{prop:alphas}, $\sum_{j=t}^i\frac{(j+1)^{2(k-1)}}{\weight_{1:j}}\leq \sum_{j=t}^i\frac{4(j+1)^4}{j^4} \leq 64i$. Hence,
% \begin{align*}
%     \sum_{i=t}^T\frac{c}{\lambda_i}\sum_{j=t}^i\frac{(j+1)^{2(k-1)}}{\weight_{1:j}}
%     &\leq \sum_{i=t}^T \frac{64}{i^{3/2}} \leq \frac{64}{t^{3/2}}+\int_t^\infty \frac{64}{x^{3/2}}\,dx = \frac{64}{t^{3/2}} + \frac{128}{t^{1/2}} \leq \frac{192}{t^{1/2}}.
% \end{align*}
% Consequently, the second term is bounded by:
% \begin{align*}
%     \sum_{t=1}^T\sum_{i=t}^T\frac{1}{\lambda_i}\sum_{j=t}^i\frac{(j+1)^{2(k-1)}}{\weight_{1:j}} 2\weight_t\|w_t\|_2^2
%     &\leq \sum_{t=1}^T \frac{384}{c}t^{5/2}\|w_t\|_2^2.
% \end{align*}
In conclusion, upon substituting \textbf{1.2.1.} and \textbf{1.2.2.} into \eqref{eq:step-1.2} and setting $c=14$, we get:
\begin{align*}
    &8C^2\sum_{t=1}^T\sqrt{\sum_{i=1}^t2H^2\|w_i-x_{i-1}\|_2^2i^4}\|w_t-x^*\|_2 \\
    \leq & AH \sum_{t=1}^T \left(\frac{5}{3c}t^{5/2}\|w_t\|_2^2 + \frac{192}{c}t^{5/2}\|w_t\|_2^2 + ct^{5/2}(\|w_t\|_2^2+\|x^*\|_2^2) \right) \\
    \leq & AH \sum_{t=1}^T 28 t^{5/2}(\|w_t\|_2^2+\|x^*\|_2^2).
\end{align*}
% The second inequality holds by choosing $c=14$.
Moreover, upon substituting \textbf{1.1.} and \textbf{1.2.} into \eqref{eq:step-1}, we get: with probability at least $1-\delta/2$,
\begin{align*}
    &\sum_{t=1}^T\|\weight_t\nabla f(x_t)-g_t\|_2\|w_t-x^*\|_2 \\
    % \leq & \sum_{t=1}^T A\left( G\sqrt{\log(2d/\delta_t)} + 14H(1+\log(2d/\delta_t)) \right) t^{5/2} (\|w_t\|_2^2+\|x^*\|_2^2).
    \leq & A\Phi\sum_{t=1}^T  Gt^{5/2}(\|w_t\|_2+\|x^*\|_2) + 28Ht^{5/2} (\|w_t\|_2^2+\|x^*\|_2^2).
\end{align*}
% \begin{align*}
%     &\sum_{t=1}^T\|\weight_t\nabla f(x_t)-g_t\|_2\|w_t-x^*\|_2 \\
%     \leq & 8C^2\sum_{t=1}^T G\sqrt{\log\frac{2d}{\delta_t}} t^{5/2}(1+\|w_t-x^*\|_2^2) + 40H\left(1+\log\frac{2d}{\delta_t}\right) t^{5/2} (\|w_t\|_2^2+\|x^*\|_2^2) \\
%     % \leq & 8C^2\sum_{t=1}^T \sqrt{\log\frac{2d}{\delta_t}}\left( Gt^{5/2} + (2G+28H)(\|w_t\|_2^2+\|x^*\|_2^2)t^{5/2} \right).
%     \leq & 8C^2G\sqrt{\log\frac{2d}{\delta_t}}T^{7/2} + 16C^2\sum_{t=1}^T\left(G\sqrt{\log\frac{2d}{\delta_t}}+20H(1+\log\frac{2d}{\delta_t})\right)t^{5/2}(\|w_t\|_2^2+\|x^*\|_2^2).
% \end{align*}

\textbf{Step 2:} We can bound $\sum_{t=1}^T\|\gamma_t\|_2\|w_t-x^*\|_2$ in a similar way. By Lemma \ref{lem:X-R-nSG} and definition of $\sigma_t$ in \eqref{eq:RDP-noise}, $R_i|R_{1:i-1}$ is mean-zero $\nSG(\sigma_{R_i})$, where
\begin{align*}
    \sigma_{R_i} 
    &= C\sqrt{d}\sigma_\calD \sigma_i
    = \frac{8\sqrt{d}\sigma_\calD C}{\rho}   \sqrt{\log_2(2T)}(G+H\max_{j\in[i]}\|w_j-x_{j-1}\|_2)i^2.
\end{align*}
Next, we can bound $\sum_{i\in I_t}\sigma_{R_i}^2$ as follows: for all $t$,
\begin{align*}
    \sum_{i\in I_t} \sigma_{R_i}^2 \geq \min_{i\in I_t}\sigma_{R_i}^2 \geq b:= \frac{64d\sigma_\calD^2C^2}{\rho^2}\log_2(2T)G^2.
\end{align*}
On the other hand, since $|I_t|\leq \log_2(2T)$,
\begin{align}
    \sum_{i\in I_t} \sigma_{R_i}^2 \leq B_t:= \frac{64d\sigma_\calD^2C^2}{\rho^2} \log_2^2(2T) (G+DH)^2 t^4.
    % \label{eq:parameter-pf-2}
    \notag
\end{align}
Hence, $B_t/b \leq \log_2(2T)\kappa^2T^4 \leq (2\kappa T)^5$ (because $\log_2(2T)\leq 2T$ and $\kappa\geq 1$). By definition of $\Phi$,
\begin{align*}
    \sqrt{\log\frac{4dT}{\delta}+\log\log\frac{B_t}{b}}
    \leq \sqrt{\log\frac{20dT\log(2\kappa T)}{\delta}} = \Phi.
\end{align*}
Recall that $A'=8\sqrt{d}\sigma_\calD C^2$. By Lemma \ref{lem:nSG-concentration}, for each $t$, with probability at least $1-\delta/2T$,
\begin{align}
    \|\gamma_t\|_2 
    &\leq C\sqrt{\sum_{i\in I_t}\sigma_{R_i}^2\left(\log\frac{4dT}{\delta}+\log\log\frac{B_t}{b}\right)}
    % \intertext{Recall that $A'=8\sqrt{dK}\sigma_\calD C^2$. Upon substituting \eqref{eq:parameter-pf-2}, we get:}
    \leq \frac{A'}{\rho}(G+DH)\Phi \log_2(2T)t^2.
    \label{eq:privacy-concentration}
\end{align}
By union bound, with probability at least $1-\delta/2$,
\begin{align*}
    \sum_{t=1}^T\|\gamma_t\|_2\|w_t-x^*\|_2 
    % &\leq \sum_{t=1}^TC\sqrt{\sum_{i\in I_t}\sigma_{R_i}^2\left(\log\frac{2dT}{\delta}+\log\log\frac{B_t}{b}\right)}\|w_t-x^*\|_2
    % &\leq \sum_{t=1}^T C\sqrt{B_t\left(\log\frac{2dT}{\delta}+\log\log\frac{B_t}{b}\right)}D
    % \intertext{Recall that $A'=8\sqrt{dK}\sigma_\calD C^2, \Phi=\sqrt{\log(16dT\log(2\kappa T)/\delta)}$, so bounding $\sum_{i\in I_t}\sigma_{R_i}^2 \leq B_t$ gives:}
    % &\leq \sum_{t=1}^T \sqrt{\log\frac{16dT\log(2\kappa T)}{\delta}}\frac{A'}{\rho} (G+DH) \log_2(2T) t^2 \|w_t-x^*\|_2 \\
    &\leq \sum_{t=1}^T \frac{A'}{\rho}(G+DH)\Phi \log_2(2T)t^2(\|w_t\|_2+\|x^*\|_2).
\end{align*}

In conclusion, we take the union bound on the results from \textbf{step 1.} and \textbf{step 2.} and substitute it back to the starting point \eqref{eq:parameter-free-base}.  Then with probability at least $1-\delta$,
\begin{align*}
    \weight_{1:T}(f(x_T)-f(x^*))
    &\leq R_T(x^*) + \sum_{t=1}^T 28AH\Phi t^{5/2}(\|w_t\|_2^2+\|x^*\|_2^2)  \\
    & \quad \ + \sum_{t=1}^T \left(AG\Phi t^{5/2} + A'(G+DH)\frac{\Phi\log_2(2T)t^2}{\rho}\right)(\|w_t\|_2+\|x^*\|_2).
\end{align*}
Define $\xi_t,\nu_t$ as in the theorem, and recall that $\weight_{1:T}\geq T^4/4$. This completes the proof.
\end{proof}

% \ThmParameterFree*

% \begin{proof}
% The regularized regret satisfies
% \begin{align*}
%     \overline \regret_T(x^*) 
%     &= \sum_{t=1}^T [\ell_t(w_t) + \xi_t\|w_t\|_2 + \nu_t\|w_t\|_2^2] - [\ell_t(x^*)+\xi_t\|x^*\|_2^2 + \nu_t\|x^*\|_2^2] \\
%     &= \regret_T(x^*) + \sum_{t=1}^T \xi_t(\|w_t\|_2-\|x^*\|_2)+\nu_t(\|w_t\|_2^2-\|x^*\|_2^2).
% \end{align*}
% Hence, upon substituting this equation into Theorem \ref{thm:parameter-free-2-norm-step-1}, we get: with probability at least $1-\delta$,
% \begin{align*}
%     f(x_T)-f(x^*) 
%     &\leq \frac{4\overline\regret_T(x^*)}{T^4} + \frac{8}{T^4} \sum_{t=1}^T \xi_t\|x^*\|_2 + \nu_t\|x^*\|_2^2 \\
%     &\leq \frac{4\regret_T(x^*)}{T^4} + \frac{8A\|x^*\|(G+28\|x^*\|H)\Phi}{\sqrt{T}} + \frac{8A'\|x^*\|(G+DH)\Phi\log_2(2T)}{\rho T}.
% \end{align*}
% The second inequality is from $\sum_{t=1}^T \xi_t \leq T\xi_t$ because $\xi_t$ is increasing with $t$ (so is $\nu_t$). 

% For the second part of the theorem, $\bar g_t = g_t+\gamma_t + \xi_t g+ 2\nu_t w$, where $g\in \partial \|w\|_2$ and thus $\|g\|_2 = 1$. Therefore,
% \begin{align*}
%     \|\bar g_t\|_2
%     &\leq \| g_t - \weight_t\nabla f(x_t)\|_2 + \|\weight_t\nabla f(x_t)\|_2 + \|\gamma_t\|_2 + \|\xi_tg\|_2 + \|2\nu_tw \|_2
% \end{align*}
% By Lipschitzness, $\|\weight_t\nabla f(x_t)\|_2 \leq Gt^3$. Since $W$ is bounded, $\|w\|_2\leq D$. For the first term, we can prove (in Equation \ref{eq:variance-concentration}) that for each $t$, with probability at least $1-\delta/2T$,
% \begin{align*}
%     \|\weight_t\nabla f(x_t) - g_t\|_2
%     &\leq 8C^2\Phi \sqrt{\sum_{i=1}^t i^4(G+H\|w_i-x_{i-1}\|_2)^2} 
%     \leq A\Phi (G+DH) t^{5/2}.
% \end{align*}
% For the third term, we can prove (in Equation \eqref{eq:privacy-concentration}) that for each $t$, with probability at least $1-\delta/2T$,
% \begin{align*}
%     \|\gamma_t\|_2 \leq \frac{A'}{\rho}(G+DH)\Phi \log_2(2T)t^2.
% \end{align*}
% In conclusion, upon taking the union bound over all $t$ and substituting the definition of $\xi_t,\nu_t$, we get: with probability at least $1-\delta$, for all $t$,
% \begin{align*}
%     \| \bar g_t \|_2 
%     &\leq A\Phi(G+DH)t^{5/2} + Gt^3 + A'\Phi(G+DH) \frac{\log_2(2T)t^2}{\rho} \\
%     &\quad \ + \left(A\Phi Gt^{5/2} +  A'\Phi(G+DH)\frac{\log_2(2T)t^2}{\rho}\right) + 56A\Phi DHt^{5/2} \\
%     &\leq Gt^3 + A(2G+57DH)\Phi t^{5/2} + 2A'(G+DH)\frac{\Phi\log_2(2T)t^2}{\rho}.
% \end{align*}
% \end{proof}


% \LemmaChangeSum*

\begin{proof}
The proof is basically re-pairing the summations:
\begin{align*}
    \sum_{i=1}^N\sum_{j=1}^i a_i b_j 
    &= a_1b_1 + a_2(b_1+b_2) + a_3(b_1+b_2+b_3) + \dotsm \\
    &= (a_1+\dotsm+a_N)b_1 + (a_2+\dotsm+a_N)b_2 + \dotsm \sum_{i=1}^T\sum_{j=0}^{T-i} a_{t-j} b_i.
\end{align*}
For the second part of the theorem, denote $B_j^i=\sum_{k=j}^i b_k$. By first part,
\begin{align*}
    LHS &= \sum_{i=1}^N a_i \sum_{j=1}^i c_j \sum_{k=j}^i b_k \\
    &= \sum_{i=1}^N\sum_{j=1}^i a_ic_j (B_j^N-B_{i+1}^N) \\
    &= \sum_{i=1}^N\sum_{j=i}^N a_jc_iB_i^N - a_jc_iB_{j+1}^N \\
    &= \sum_{i=1}^N\sum_{j=i}^Na_jc_iB_i^j.
\end{align*}
We then recover the lemma once we substitute $B_i^j=\sum_{k=i}^j b_k$.
\end{proof}

% \begin{lemma}
% \label{lem:sub-gaussian-sum}
% Suppose $\P[\|R_t\|_*\geq \epsilon]\leq 2\exp(-\epsilon^2/2\sigma_t^2)$ and $\sigma_t^2\leq \sigma_{t+1}^2$ for all $t$. Then with probability $\geq 1-\delta$, for all $t$,
% \begin{align*}
%     \sum_{i\in I_t}\|R_i\|_*^2 \leq 2\log(2t)\sigma^2 \log\left(\frac{2T\log(2t)}{\delta}\right).
% \end{align*}
% \end{lemma}

% \begin{proof}
% By pigeonhole principle and union bound, for each $t$,
% \begin{align*}
%     \P\left[\sum_{i\in I_t} \|R_i\|_*^2 \geq \epsilon \right] 
%     &\leq \P\left[ \exists i\in I_t,\, \|R_i\|_*^2 \geq \frac{\epsilon}{|I_t|} \right] \\
%     &\leq \sum_{[a,b]\in I_t} \P\left[ \|r_{[a:b]}\|_*^2 \geq \frac{\epsilon}{|I|} \right] \\
%     &\leq 2\log(2t) \exp\left(\frac{-\epsilon}{2\log(2t)\sigma^2}\right).
% \end{align*}
% Equivalently, for each $t$, with probability at least $1-\delta$, 
% \begin{align*}
%     \sum_{[a:b]\in I_t}\|r_{[a:b]}\|_*^2 \leq 2\log(2t)\sigma^2 \log\left(\frac{2\log(2t)}{\delta}\right).
% \end{align*}
% After replacing $\delta$ with $\delta/T$ and taking the union bound over all $t$, with probability at least $1-\delta$,
% \begin{align*}
%     \forall t, \, \sum_{[a:b]\in I_t}\|r_{[a:b]}\|_*^2 \leq 2\log(2t)\sigma^2 \log\left(\frac{2T\log(2t)}{\delta}\right).
% \end{align*}
% \end{proof}